\documentclass[11pt]{article}
\usepackage{amsmath,amssymb,amsthm,hyperref}
\newtheorem{lemma}{Lemma}
\title{Switches between reduced errors $-1$ and $-3$}
\author{Finite counterexample exploration for Erd\H{o}s 243; no priority claim}
\date{6 November 2026}
\begin{document}
\maketitle

\noindent\textbf{Status.} No infinite counterexample is obtained.
The entire reduced-error alphabet $\{-1,-3\}$ has not been excluded.
We prove necessary conditions on its switches and report an exact finite
search of 140,000 rational starts. All those starts reached zero error.
Neither the finite search nor the local conditions settle Erd\H{o}s 243.

\paragraph{The exact states.}
For the reduced rational tail $r_n=u_n/v_n$, let
\[
 e_n=v_n-(a_n-1)u_n,\quad w_n=u_n-e_n,\quad
 h_n=\gcd(w_n,a_nv_n),\qquad
 u_{n+1}=w_n/h_n,\quad v_{n+1}=a_nv_n/h_n.
\]
All states and multipliers are positive integers, and
$\gcd(u_n,v_n)=1$. In the finite search, the multiplier is chosen by
$a_n=1+\operatorname{round}(v_n/u_n)$, with half-integers rounded
upward, following
\href{https://math.colgate.edu/~integers/aa28/aa28.pdf}{Koizumi's
pseudo-greedy recurrence}. This is also the eventual choice under the
original near-squaring hypotheses.

Define $\Delta_n=a_{n+1}-(a_n^2-a_n+1)$. Direct substitution gives
\begin{equation}\label{eq:error}
 h_ne_{n+1}=a_n^2e_n-\Delta_n(u_n-e_n).
\end{equation}
Indeed, using $v_n=a_nu_n-w_n$, the left side equals
$a_nv_n-(a_{n+1}-1)w_n
=a_n^2u_n-(a_n+a_{n+1}-1)w_n$.

\begin{lemma}[Necessary conditions on an uncancelled switch]
Suppose $h_n=1$ and $e_n,e_{n+1}$ are distinct elements of
$\{-1,-3\}$. Then $\gcd(u_{n+1},3)=1$ and $3$ is a square
modulo $u_{n+1}$. In particular, every odd prime divisor $\ell$
of $u_{n+1}$ satisfies $\ell\equiv1$ or $11\pmod{12}$,
and the power of $2$ dividing $u_{n+1}$ is at most $2$.

If $u_{n+1}$ is even, then it is $2$ modulo $4$, $a_{n+1}$ is even,
$h_{n+1}$ is odd, and $v_{n+2}$ is even. If also $e_{n+2}$ is odd,
then $h_{n+2}$ is even and hence at least $2$.
\end{lemma}
\begin{proof}
Since $h_n=1$, equation~\eqref{eq:error} gives
$e_{n+1}\equiv a_n^2e_n\pmod{u_{n+1}}$.
Also $\gcd(u_{n+1},e_n)=1$, because $u_{n+1}=u_n-e_n$ and
$\gcd(u_n,e_n)=1$. Primitive coprimality similarly gives
$\gcd(u_{n+1},e_{n+1})=1$. Thus $3$ is a unit modulo $u_{n+1}$.
For the switch $-1\to-3$ one has $a_n^2\equiv3$;
for $-3\to-1$, one has $3a_n^2\equiv1$, hence
$(3a_n)^2\equiv3$. Quadratic reciprocity gives the stated odd-prime
classes: $(3/\ell)=(\ell/3)(-1)^{(\ell-1)/2}$.
There is no square congruent to $3$ modulo $4$.

If $u_{n+1}$ is even, $v_{n+1}=a_nv_n$ is odd by coprimality,
so $a_n$ and $v_n$ are odd. Hence
$a_n^2e_n-e_{n+1}\equiv e_n-e_{n+1}\equiv\pm2\pmod8$.
Its $2$-adic valuation is exactly one. Division by the even
$u_{n+1}$ in~\eqref{eq:error} shows that $u_{n+1}\equiv2\pmod4$
and $\Delta_n$ is odd. Since $a_n^2-a_n+1$ is odd,
$a_{n+1}$ is even. Now
$w_{n+1}=u_{n+1}-e_{n+1}$ is odd, so $h_{n+1}$ is odd and
$v_{n+2}=a_{n+1}v_{n+1}/h_{n+1}$ is even.
Thus $u_{n+2}$ is odd. If $e_{n+2}$ is odd, the defining error
identity makes $a_{n+2}$ even. Both $w_{n+2}$ and
$a_{n+2}v_{n+2}$ are even, proving $2\mid h_{n+2}$.
\end{proof}

The hypothesis on the third error matters. For instance,
$(u_0,v_0)=(21,83)$ gives $e_0=-1$, then
$(u_1,v_1)=(22,415)$ and $e_1=-3$, followed by a cancellation
$h_1=25$ and $(u_2,v_2)=(1,332)$ with $e_2=0$. This is a
finite prefix followed by a Sylvester tail, not a counterexample.

\paragraph{A family with an enforced first switch.}
For each
\[
 p\in\{11,13,23,37,47,61,73\},\quad b\in\{1,2\},
\]
let $r$ range over the two odd representatives with
$3\le r<2p+3$ and $r^2\equiv3\pmod p$. For
$0\le t<5000$, set
\[
 u_0=bp-1,\quad a_0=r+2pt,\quad
 v_0=(a_0-1)u_0-1.
\]
These are positive primitive rational states. Rounding gives $e_0=-1$,
and $w_0=bp$. The multiplier is odd and not divisible by $p$, so
$h_0=\gcd(bp,a_0^2)=1$. Modulo $bp$, the next denominator is
$-a_0^2\equiv-3$. Since $bp\ge11$, its centered error is exactly
$-3$ and the next state is still primitive. Thus all
$7\cdot2\cdot2\cdot5000=140000$ starts have the first two
reduced errors $-1,-3$.

The standalone executable \texttt{erdos243-small-alphabet-search.mjs}
first follows each start exactly while its errors stay in this alphabet.
It then follows \emph{every} start to an exact zero error using Koizumi's
modular recomputation. To obtain $D_n\bmod C_n$, the algorithm starts
modulo $\prod_{j\le n}C_j$, replays each preceding division by $C_j$
while deleting that factor from the modulus, and ends modulo $C_n$.
No pairwise coprimality of the raw numerators is assumed. Its gcd gives
the reduced numerator and error. The computation has explicit caps;
none was reached in this experiment.

All 140,000 starts reached zero, with at most $22$ recorded states
including the first zero state. Thus each has a Sylvester tail.
The longest initial stay in $\{-1,-3\}$ had $13$ errors and started at
\[
 (u_0,v_0)=(10,582499).
\]
Its reduced numerator/error/cancellation triples were
\[
\begin{array}{c|c|c}
u_n&e_n&h_n\\\hline
10&-1&1\\
11,14,17,20,23,26,29,32,35,38&-3&1\\
41&-3&11\\
4&-1&1\\
5&1&1\\
4&1&3\\
1&0&\text{terminal}
\end{array}
\]
The largest multiplier in that alphabet prefix had $63,953$ bits;
only its bit length is stored in the compact report, since the seed
determines it exactly. The longest full continuation instead began at
$(121,45253757)$ and reached its first zero at its $22$nd state.
The finite output and seed-class champions are retained in
\texttt{erdos243-small-alphabet-results.json}.

Direct reduced arithmetic was compared with independent unreduced
updates for every reduced $0<p\le q\le40$ up to six steps, and the
modular continuation was compared with every step of all $28$ class
champions' alphabet prefixes. The known $11/29$ trajectory also checks
the raw-to-reduced gcd handling. Each observed uncancelled switch
was checked against the lemma, including the forced even cancellation
when the subsequent error remained odd.

\paragraph{Why a single-prime barrier is insufficient here.}
An abstract increasing walk can use increments $1$ and $3$, obey
$\gcd(u_n,u_{n+1})=1$, and avoid multiples of $5$ indefinitely.
For example, repeat the following residue pattern with a shift of $30$:
\[
 1,2,3,4,7,8,11,12,13,16,17,18,19,22,23,26,27,28,31.
\]
Every jump of length $3$ starts at a number coprime to $3$; the
shift of $30$ preserves this and avoidance of $5$.
This is only a scalar walk. No positive denominator sequence is
specified, and it is not an orbit or a disproof. It shows that the
fixed-magnitude lattice barrier cannot be extended to this alphabet
using only small increments and adjacent coprimality. The exact
denominator and switch constraints must do further work.

\paragraph{Unresolved construction obligation.}
A full counterexample still needs one fixed positive rational start,
exact integral updates at every index, $u_n\to\infty$, and
$e_n/u_n\to0$ with errors nonzero arbitrarily late. These finite
trajectories and necessary switch conditions establish none of that
infinite existence statement. No exclusion of all such infinite orbits
has been proved either.
\end{document}
